\section{Methodology}

\subsection{Data Source and Sample}

The analysis uses a subset of the United States Department of Education College Scorecard, a public database that reports institutional characteristics, costs, student outcomes, and debt for postsecondary institutions across the country. The Scorecard was developed to help students and families make informed decisions about higher education by providing transparent and comparable data on costs and expected earnings.

The data are collected through federal reporting requirements. Institutions that participate in Title IV federal student aid programs must report enrollment, completion, tuition, and financial aid data to the Department of Education. Earnings outcomes are drawn from de-identified tax records matched to former students through the Treasury Department, which provides a reliable measure of labor-market returns. The collection process imposes standardized definitions across institutions, which strengthens internal comparability. However, the earnings figures are limited to students who received federal financial aid and whose tax records could be matched. This means the earnings data may not represent all graduates and may understate outcomes for students from higher-income families who did not receive aid.

The dataset contains 6,273 rows and 24 columns. Each row represents one postsecondary institution identified by a unique unit identifier. The columns capture institutional identity, admissions selectivity, demographic designations, tuition, student debt, and post-graduation earnings. The key variables for this study are in-state tuition (TUITIONFEE\_IN), median earnings ten years after entry (MD\_EARN\_WNE\_P10), institutional control type (CONTROL), median graduate debt (GRAD\_DEBT\_MDN), and undergraduate enrollment (UGDS). The full list of variables is available in the College Scorecard data dictionary.

\subsection{Data Cleaning}

The raw dataset required substantial cleaning before analysis. Five variables were selected as central to the research question: in-state tuition, median post-graduation earnings, median graduate debt, institutional control type, and undergraduate enrollment. Two variables were excluded from the primary analysis. SAT average scores were missing for 83.5\% of institutions and admission rates for 69.6\%. Imputing values for the majority of observations would introduce bias and undermine the validity of downstream inferences. Institutions reporting zero in-state tuition were also removed, because the investment-return framing of this study requires a positive cost baseline against which to measure earnings.

The cleaning pipeline addressed six issues sequentially. First, the graduate debt column contained ``PrivacySuppressed'' string values for 1,228 institutions. These were coerced to numeric NaN using \texttt{pd.to\_numeric} with \texttt{errors='coerce'}. Second, rows missing the target variable (MD\_EARN\_WNE\_P10) or the primary predictor (TUITIONFEE\_IN) were dropped. Third, rows missing graduate debt after the privacy suppression conversion were dropped. Fourth, missing undergraduate enrollment values were imputed with the median enrollment within each sector rather than a global median, because enrollment distributions differ substantially across public, private nonprofit, and for-profit institutions. Fifth, duplicate UNITID values were detected and removed. Sixth, outliers were trimmed using the interquartile range method. Earnings outliers were removed symmetrically across the full sample. Tuition outliers were removed within each sector separately, because tuition distributions differ dramatically by control type and a global IQR filter would disproportionately remove legitimate observations from the private nonprofit tier.

The final cleaning step created a human-readable sector label from the numeric CONTROL code, mapping 1 to Public, 2 to Private Nonprofit, and 3 to For-Profit. The pipeline reduced the dataset from 6,273 to 2,894 observations. The largest reduction came from the 2,915 institutions that lacked valid tuition or earnings data. The median imputation for undergraduate enrollment retained 257 institutions that would otherwise have been dropped, and the sector-specific IQR outlier removal stripped 167 extreme values. The final dataset contains no missing values in any of the five target variables.

\subsection{Descriptive Framework}

Summary statistics were computed stratified by institutional control type. Panel A reports the sector composition of the cleaned dataset. The final sample of 2,894 institutions comprises 1,430 public institutions (49.4%), 1,104 private nonprofit institutions (38.1%), and 360 for-profit institutions (12.4%). Panel B reports the mean, median, standard deviation, and interquartile range for in-state tuition, median post-graduation earnings, median graduate debt, and undergraduate enrollment, computed separately for each sector and for the full sample.

The descriptive statistics serve two purposes. They establish the central tendencies and dispersions that any model of the tuition-earnings relationship must account for. They also reveal the sector-level disparities that motivate the stratified modeling approach. If earnings and tuition differ systematically by sector, a model that pools all institutions into a single slope will misrepresent the relationship for every sector.

\subsection{Exploratory Data Analysis}

Six exploratory questions were posed, moving from univariate description to joint and contextual relationships. The first question examined the distribution of post-graduation earnings across sectors using grouped summary statistics, kernel density estimation, and boxplots. The second examined in-state tuition distributions using grouped statistics and violin plots, with the national median as a benchmark. The third assessed the joint tuition-earnings relationship through Spearman rank-order correlations and scatter plots with sector-stratified regression lines. The fourth examined graduate debt as a mediating variable, computing Spearman correlations between tuition and debt alongside median debt-to-earnings ratios for each sector. The fifth tested whether the tuition-earnings link persisted within admission selectivity tiers, stratifying institutions into highly selective (admission rate below 40\%), moderately selective (40\% to 70\%), and less selective (above 70\%) categories. The sixth compared minority-serving institutions against non-minority-serving institutions within each sector using grouped means and medians for tuition, earnings, and debt.

These exploratory analyses established three patterns that directly informed the modeling strategy. Earnings and tuition differ systematically by sector. The within-sector tuition-earnings correlation is strongest among public institutions and weakest among for-profits. And the debt-to-earnings ratio is substantially lower for public graduates than for graduates of private nonprofit or for-profit institutions.

\subsection{Statistical Inference Framework}

Before estimating the regression models, the distribution of median earnings was tested for normality using the D'Agostino K-squared test. The null hypothesis was that the earnings data follow a normal distribution. The alternative was that they deviate from normality. Because the earnings distribution was expected to show positive skew and multimodality given the sector-level differences observed in the exploratory analysis, a significant test result would justify the use of Spearman's rank correlation for the preliminary association test. The significance level was set at $\alpha = 0.05$.

The D'Agostino K-squared test produces a test statistic and p-value, and if the data deviate significantly from normality, the analysis uses Spearman's rank correlation rather than Pearson's r. Visual diagnostics including a histogram and Q-Q plot supplement the formal test to assess whether deviations from normality are practically severe or merely a consequence of the large sample size. The Spearman correlation coefficient and its associated p-value are reported in the Results section.

\subsection{Regression Modeling Strategy}

The research question asks whether a significant relationship exists between in-state tuition and median post-graduation earnings and whether this relationship differs across public, private nonprofit, and for-profit institutions. A nested sequence of three ordinary least squares models was specified to answer both parts of the question. This design follows the institution-level higher-education research literature. \textcite[pp.~41--45]{faber2017} and \textcite[pp.~687--700]{faber2021} model earnings using College Scorecard institutional attributes including state, enrollment, and tuition. \textcite[pp.~36--40]{muse2024} likewise model College Scorecard earnings with tuition, institutional type, size, degree offerings, geography, and selectivity controls.

All three models use the natural logarithm of in-state tuition and the natural logarithm of median post-graduation earnings. Both variables are positive and right-skewed, and the log transformation reduces skewness while producing coefficients that are interpretable as elasticities. A coefficient of 0.10 means that a 1\% increase in tuition is associated with approximately 0.10\% higher earnings. \textcite[pp.~37, 40]{muse2024} apply the same logarithmic transformation to College Scorecard median earnings. Undergraduate enrollment enters as $\log(1 + \text{UGDS})$ because one institution reports zero enrollment, and the log of zero is undefined.

All models use HC3 heteroskedasticity-consistent standard errors. \textcite[pp.~307--309]{mackinnon1985} derive the jackknife-based HC3 covariance estimator for regression inference under heteroskedasticity, and the Breusch-Pagan test for constant residual variance is reported for each model. The estimates describe conditional associations and must not be interpreted as causal effects of tuition on earnings.

\subsubsection{Model 1: Unconditional Association}

Model 1 estimates the overall association between in-state tuition and median earnings without adjusting for institutional characteristics. The specification is

\[
\log(\text{Earnings})_i = \beta_0 + \beta_1 \log(\text{Tuition})_i + \varepsilon_i
\]

where $i$ indexes institutions. The coefficient $\beta_1$ represents the pooled tuition elasticity across all sectors. This model addresses the first part of the research question but pools public, private nonprofit, and for-profit institutions into a single slope, which the exploratory analysis has shown to be misleading.

\subsubsection{Model 2: Sector Interaction Model}

Model 2 adds institutional sector indicators and interacts each sector with log tuition. The specification is

\[
\begin{aligned}
\log(\text{Earnings})_i = &\ \beta_0 + \beta_1 \log(\text{Tuition})_i \\
&+ \beta_2 \text{Nonprofit}_i + \beta_3 \text{ForProfit}_i \\
&+ \beta_4 \log(\text{Tuition})_i \times \text{Nonprofit}_i \\
&+ \beta_5 \log(\text{Tuition})_i \times \text{ForProfit}_i + \varepsilon_i
\end{aligned}
\]

Public institutions are the reference category. The coefficients $\beta_4$ and $\beta_5$ test whether the tuition elasticity for private nonprofit and for-profit institutions differs from the public institution elasticity. A joint HC3 Wald test evaluates whether the two interaction coefficients are jointly zero, which directly answers the second part of the research question. \textcite[pp.~36, 39]{muse2024} distinguish the same three control types when modeling College Scorecard earnings, and \textcite[pp.~687--700]{faber2021} show that institutional attributes change the interpretation of tuition-earnings relationships.

\subsubsection{Model 3: Institutional-Control Robustness Model}

Model 3 retains the Model 2 interaction terms and adds three prespecified blocks of institutional controls. The specification is

\[
\begin{aligned}
\log(\text{Earnings})_i = &\ \beta_0 + \beta_1 \log(\text{Tuition})_i \\
&+ \beta_2 \text{Nonprofit}_i + \beta_3 \text{ForProfit}_i \\
&+ \beta_4 \log(\text{Tuition})_i \times \text{Nonprofit}_i \\
&+ \beta_5 \log(\text{Tuition})_i \times \text{ForProfit}_i \\
&+ \beta_6 \log(1 + \text{Enrollment})_i \\
&+ \gamma \cdot \text{Degree}_i + \delta \cdot \text{State}_i + \varepsilon_i
\end{aligned}
\]

Undergraduate enrollment enters as $\log(1 + \text{UGDS})$. Predominant degree enters categorically with bachelor's degree as the reference category, using the PREDDEG variable where 0 indicates non-degree-granting, 1 indicates certificate or associate, 2 indicates associate, 3 indicates bachelor's, and 4 indicates graduate-focused. Institution state enters as fixed effects using the STABBR variable, with Guam, Puerto Rico, and the U.S. Virgin Islands grouped as ``US Territory'' to avoid separate fixed effects supported by only one or two observations.

The control set follows the closest published precedents that can be mapped to this dataset. \textcite[pp.~42--45]{faber2017} includes state, enrollment, and tuition among blocked College Scorecard institutional attributes before estimating earnings with OLS. \textcite[pp.~36--40]{muse2024} include selectivity, region, institutional type, size, degree offerings, tuition, and log earnings. Graduate debt is excluded from the primary model because it may lie on the causal pathway from tuition to later earnings outcomes. Admission rate is handled separately in a disclosed sensitivity analysis because it is observed for only 50.7\% of the cleaned sample.

The robustness question is whether the sector-specific tuition elasticities remain positive and whether the interaction terms remain jointly significant after these controls. A joint HC3 Wald test evaluates whether the additional institutional controls contribute explanatory information beyond Model 2.

\subsubsection{Admission-Rate Sensitivity Analysis}

Admission rate is added only as a disclosed sensitivity analysis and does not replace the full-sample Model 3. \textcite[pp.~35--37]{muse2024} treat admission rate as a selectivity indicator but document that selectivity measures are missing for many institutions, and the same problem is material here. Only 1,466 of 2,894 institutions report ADM\_RATE, and the reporting subset contains 908 private nonprofits, 503 public institutions, and only 55 for-profits.

Two models are estimated on exactly that restricted sample. The first repeats the Model 3 specification without admission rate. The second adds admission rate as a linear predictor. Comparing the two separates the effect of adding selectivity from the effect of changing the estimation sample. Missing admission rates are unlikely to be random, because selective institutions are disproportionately likely to report their rates. The restricted results must not be generalized to all institutions.

\subsection{Hypothesis Testing and Inference}

All hypothesis tests use a significance level of $\alpha = 0.05$. The null hypothesis for the overall tuition-earnings association is $\text{H}_0: \beta_1 = 0$ against $\text{H}_1: \beta_1 \neq 0$. The null hypothesis for the sector differences is that the tuition-by-sector interaction coefficients are jointly zero, tested using an HC3 Wald statistic.

Model fit is assessed using adjusted R-squared, log RMSE, AIC, and BIC. The Breusch-Pagan test for heteroskedasticity is reported for each model to validate the use of HC3 standard errors. The sector-specific tuition elasticities and their 95\% confidence intervals are computed using the delta method via linear hypothesis tests on the combined coefficients. For each sector, the elasticity is the sum of $\beta_1$ and the relevant interaction term, and the standard error is computed from the HC3 covariance matrix.

\subsection{Software}

All analyses were conducted in Python 3.12. Data manipulation used pandas (version 3.0) and numpy (version 2.5). Normality testing and correlation analysis used SciPy (version 1.18). Regression models were estimated using statsmodels (version 0.14) with the \texttt{smf.ols} interface and HC3 covariance specification. Visualizations were generated using matplotlib (version 3.11) and seaborn (version 0.13), styled following the Economist-aligned Quantitative Visualization and Reporting Standard. The complete analysis code is available in the companion Jupyter notebook.
